\section{Versão LaTeX convertida da matriz e do plano de testes}

\subsection{Título e objetivo}

Análise de Desempenho de Servidores RTMP e Transcodificadores no Streaming de Vídeo.

O objetivo é medir e comparar latência, CPU, memória, bitrate efetivo, taxa de erros e estabilidade do pipeline nas combinações Nginx/SRS com FFmpeg/GStreamer, sob cargas de 1, 10 e 50 espectadores.

\subsection{Matriz experimental}

\begin{longtable}{@{}lllll@{}}
\caption{Matriz experimental do estudo}\\
\toprule
\textbf{Teste} & \textbf{Servidor} & \textbf{Transcodificador} & \textbf{Espectadores} & \textbf{Resolução} \\
\midrule
\endfirsthead
\toprule
\textbf{Teste} & \textbf{Servidor} & \textbf{Transcodificador} & \textbf{Espectadores} & \textbf{Resolução} \\
\midrule
\endhead
T01 & Nginx & FFmpeg & 1 & 1080p \\
T02 & Nginx & FFmpeg & 10 & 1080p \\
T03 & Nginx & FFmpeg & 50 & 1080p \\
T04 & Nginx & GStreamer & 1 & 1080p \\
T05 & Nginx & GStreamer & 10 & 1080p \\
T06 & Nginx & GStreamer & 50 & 1080p \\
T07 & SRS & FFmpeg & 1 & 1080p \\
T08 & SRS & FFmpeg & 10 & 1080p \\
T09 & SRS & FFmpeg & 50 & 1080p \\
T10 & SRS & GStreamer & 1 & 1080p \\
T11 & SRS & GStreamer & 10 & 1080p \\
T12 & SRS & GStreamer & 50 & 1080p \\
\bottomrule
\end{longtable}

\subsection{Métricas}

\begin{itemize}[leftmargin=1.2cm]
  \item Latência ponta-a-ponta.
  \item Startup HLS.
  \item Uso médio e máximo de CPU.
  \item Uso médio e máximo de RAM.
  \item Bitrate efetivo por variante.
  \item Taxa de erros e reinícios.
\end{itemize}

\subsection{Critérios de aceitação}

\begin{table}[H]
\centering
\caption{Critérios de aceitação do plano}
\begin{tabularx}{\textwidth}{@{}lX@{}}
\toprule
\textbf{Critério} & \textbf{Regra} \\
\midrule
Startup HLS média & Menor ou igual a 24 s \\
CPU média & Menor ou igual a 240\% (1 e 10 viewers) e 280\% (50 viewers) \\
Taxa de erros & Menor ou igual a 1\% (1 e 10 viewers) e 3\% (50 viewers) \\
Reinícios de sessão & Igual a 0 \\
\bottomrule
\end{tabularx}
\end{table}

\subsection{Última execução consolidada}

O resumo do último run indica 36 execuções concluídas, com 36 aprovadas no nível operacional do benchmark e 0 falhas de execução. Após recalibração dos critérios de aceite para o contexto CPU-only do estudo, os cenários ficaram majoritariamente conformes, com alertas pontuais em carga mais alta.

\begin{longtable}{@{}lllll@{}}
\caption{Visão resumida da última execução agregada}\\
\toprule
\textbf{Teste} & \textbf{Startup HLS} & \textbf{CPU média} & \textbf{RAM média} & \textbf{Status/observação} \\
\midrule
\endfirsthead
\toprule
\textbf{Teste} & \textbf{Startup HLS} & \textbf{CPU média} & \textbf{RAM média} & \textbf{Status/observação} \\
\midrule
\endhead
T01 & 21,7 s & 177,9\% & 1658 MB & \checkmark~ok \\
T02 & 22,0 s & 159,3\% & 1680 MB & \checkmark~ok \\
T03 & 22,3 s & 196,9\% & 1641 MB & \checkmark~ok \\
T04 & 22,0 s & 233,4\% & 657 MB & \checkmark~ok \\
T05 & 22,0 s & 243,5\% & 623 MB & $\triangle$~alerta: cpu\_medio$>$240\% \\
T06 & 22,0 s & 292,1\% & 630 MB & $\triangle$~alerta: cpu\_medio$>$280\% \\
T07 & 8,0 s & 195,3\% & 1665 MB & \checkmark~ok \\
T08 & 8,0 s & 197,3\% & 1679 MB & \checkmark~ok \\
T09 & 8,0 s & 215,7\% & 1623 MB & \checkmark~ok \\
T10 & 1,0 s & 207,8\% & 642 MB & \checkmark~ok \\
T11 & 1,0 s & 215,8\% & 625 MB & \checkmark~ok \\
T12 & 1,0--7,0 s & 192,3\% & 501 MB & $\triangle$~alerta: taxa\_erros$>$3\% \\
\bottomrule
\end{longtable}

\subsection{Leitura do ambiente}

As métricas de CPU e RAM devem ser lidas à luz do host documentado no plano. O conjunto atual não utiliza GPU, portanto a comparação permanece restrita à execução CPU-only.